\subsubsection{Statement B.2r: geometric torus exclusion and bounded exponent}
\begin{proposition}\label{B:B2-no-torus-exponent}
Let $k$ be a field of characteristic $p>0$.
\begin{enumerate}
\item If $Y$ is a proper normal $k$-variety, then
$(\operatorname{Pic}^{\tau}_{Y/k})_{\bar k}$ contains no nontrivial
torus.  Geometric normality of $Y$ is unnecessary.
\item Let $H$ be a commutative $k$-group scheme of finite type such that
$H_{\bar k}$ contains no nontrivial torus.  Then $H$ has a largest
abelian subvariety $A_H$, and $H/A_H$ is annihilated by one positive
integer.  More precisely, for all sufficiently large $n$,
\[
 A_H=\operatorname{im}([p^n]:H^0\longrightarrow H^0).
\]
\item If $l/k$ is any finite field extension and $P$ is a
proper commutative $l$-group scheme of finite type, then
$\operatorname{Res}_{l/k}P$ contains no geometric torus.  Thus (ii)
applies to every closed subgroup scheme of this restriction of scalars.
\end{enumerate}
\end{proposition}
\begin{proof}
We first prove (i), including the assertion on arbitrary test schemes
needed for a statement about group schemes.  Put $\Omega=\bar k$,
$X=Y_\Omega$, and let
\[
 \nu:Z=(X_{\mathrm{red}})^{\mathrm{norm}}\longrightarrow X.
\]
Tanaka's Lemma~2.2(2), including its proof on pp.~6--8, says that the
normalization of the reduction of any field extension of a normal
variety is a universal homeomorphism
\cite[Lemma 2.2(2)]{Tanaka2015D}.
Here normalization is finite, since $X$ is of finite type over a field.
Consequently $\nu$ is a finite universal homeomorphism.  Its kernel on
structure sheaves is nilpotent.  Kelly's Corollary~20 therefore supplies
an integer $e$ and an absolute Frobenius factorization
\[
 \sigma:X\longrightarrow Z,\qquad
 \nu\sigma=F_X^e,\qquad \sigma\nu=F_Z^e
\]
\cite[Corollary 20, p.~9]{Kelly2018D}.
The morphism $\sigma$ is $F_\Omega^e$-semilinear; it need not be an
$\Omega$-morphism.

We claim that the kernel of
$\nu^*:\operatorname{Pic}_{X/\Omega}\to
\operatorname{Pic}_{Z/\Omega}$ is annihilated by $p^e$ as an fppf
sheaf.  Indeed, after an fppf localization on a test scheme $T$, an
element of this kernel is represented by a line bundle $L$ on $X_T$
such that $\nu_T^*L$ is pulled back from $T$.  The two semilinear maps
$\sigma$ and $F_T^e$ define a morphism
$\sigma\times F_T^e:X_T\to Z_T$.  Their composite with $\nu_T$ is the
absolute Frobenius $F_{X_T}^e$.  Pulling back $\nu_T^*L$ along this
morphism shows that $L^{\otimes p^e}$ is pulled back from $T$.
This proves the claim uniformly for all $T$; no identification of a
relative Frobenius with multiplication on a Picard scheme is used.

The normal proper $\Omega$-scheme $Z$ has proper torsion Picard
component, by Kleiman's Theorem~5.4 and Remarks~5.6 and~6.19
\cite{Kleiman2005}.  A homomorphism from a torus to a proper algebraic
group is zero: its scheme-theoretic image is a torus quotient and is
also proper, hence trivial.  Therefore any torus in
$\operatorname{Pic}^{\tau}_{X/\Omega}$ lies in the kernel of $\nu^*$.
Multiplication by $p^e$ is surjective on a torus, whereas it is zero on
this kernel.  The torus is consequently trivial.  Picard base change
identifies the group in (i) with $\operatorname{Pic}^{\tau}_{X/\Omega}$.

For (ii), take the closed scheme-theoretic images
\[
 H_n=\operatorname{im}([p^n]:H^0\longrightarrow H^0).
\]
The homomorphism theorem makes each map onto its image faithfully
flat \cite[Theorem 5.39, p.~108]{MilneAG2022B}.  The $H_n$ form a
descending chain of closed subgroup schemes in the noetherian scheme
$H^0$, so their defining ideal sheaves form an ascending chain and
stabilize.  Write $K=H_N=H_{N+1}=\cdots$.  Thus
$[p]:K\to K$ is faithfully flat.  Scheme-theoretic images commute with
flat field extension, so these assertions remain true over $\Omega$.

Over $\Omega$, the reduction $K_{\Omega,\mathrm{red}}$ is a smooth
connected subgroup.  The quotient
$Q=K_\Omega/K_{\Omega,\mathrm{red}}$ is finite and infinitesimal.
It is annihilated by a power of $p$.  For completeness, its local
augmentation ideal $I$ is nilpotent, and $[p]^*I\subset I^2$, since
$d[p]=p\,\mathrm{id}=0$ on the cotangent space.  Iteration sends $I$
into $I^{2^r}$ and is zero for large $r$.
But $[p]$ is faithfully flat on $Q$, so a power of it can be zero only
when $Q=0$.  Hence $K_\Omega$ is smooth.

Apply the Barsotti--Chevalley theorem over $\Omega$
\cite[Theorem 8.27, pp.~154--155]{MilneAG2022B}:
\[
 0\longrightarrow U\longrightarrow K_\Omega\longrightarrow B
 \longrightarrow0,
\]
where $U$ is smooth connected affine and $B$ is an abelian variety.
The group $U$ has no torus, so it is unipotent.  Some $[p^a]$ kills
$U$: embed it in a group of upper triangular unipotent matrices and
choose $p^a$ at least the matrix size.  Therefore $[p^a]$ on $K_\Omega$
factors through $B$.  It is surjective, so $K_\Omega$ is proper.
It follows that $K$ is a smooth proper connected $k$-group, hence an
abelian variety.  This deduction is made on the subgroup $K$ already
defined over $k$; it does not assume that a maximal abelian subvariety
commutes with an imperfect field extension.

Any abelian subvariety $A\subset H$ lies in $H^0$ and satisfies
$[p^N]A=A$, whence $A\subset K$.  Thus $K=A_H$ is the largest such
subvariety.  Finally $H/H^0$ is a finite etale group scheme.  If $c$
annihilates it, then $[cp^N]H\subset A_H$, proving the bounded-exponent
assertion including all infinitesimal structure and components.

For (iii), restriction of scalars exists as a finite-type algebraic
group: a proper algebraic group is quasi-projective, and restriction
along a finite field extension is representable for quasi-projective
schemes.  This is the restriction-of-scalars construction reviewed in
\cite[Section 2i, pp.~56--59]{MilneAG2022B}.
Over $\Omega$, write
$l\otimes_k\Omega=\prod_\sigma A_\sigma$, where each $A_\sigma$ is
local Artinian with residue field $\Omega$.  Restriction of scalars
commutes with this product and with field extension.  For each factor,
write $J_\sigma$ for its nilpotent augmentation ideal; choose a common
$r$ with $J_\sigma^r=0$ for all $\sigma$.  Reduction gives a homomorphism
\[
 \epsilon:(\operatorname{Res}_{l/k}P)_\Omega
       \longrightarrow\prod_\sigma P\otimes_{l,\sigma}\Omega.
\]
We show that $\ker\epsilon$ is killed by a fixed power of $p$.
We work on one factor and write $A=A_\sigma$ and $J=J_\sigma$.
For any $\Omega$-algebra $R$, kernel points on this factor are maps
$\operatorname{Spec}(R\otimes_\Omega A)\to P_A$ reducing to the identity
modulo $R\otimes_\Omega J$.  Such a map has image in an affine
neighborhood of the identity.  On the completed local algebra along
the identity section, $[p]^*$ sends the augmentation ideal into its
square, again because $d[p]=0$.  Therefore $[p]^e$ sends an identity
congruence modulo $R\otimes J$ to one modulo $R\otimes J^{2^e}$.
Choose $2^e\ge r$.  The same $[p^e]$ then kills the kernel on every
test algebra $R$, hence as a group scheme.  This argument does not
require $P$ to be smooth or the evaluation map to be surjective.

The same exponent works on all finitely many factors.  Any geometric
torus in the restriction of scalars maps trivially to the proper
product $\prod_\sigma P\otimes_{l,\sigma}\Omega$.  It is therefore contained in this
bounded-$p$-exponent kernel, and so is trivial.  A closed subgroup
inherits the same exclusion, completing the proof.
\end{proof}

\begin{remark}
This proposition is an auxiliary result for the weak Picard
stabilization argument.  It does not by itself prove that images of
Picard schemes on the permitted model category stabilize.  Its
bounded-exponent assertion is a statement about group schemes, not
only their geometric points.
\end{remark}
